\section{Discussion}

The practical contribution of projection is its guarantee rather than a
universal field-score gain: as a retraining-free layer over compatible
field-wise probabilities, it enforces the supplied hierarchy in every arm and
improves exact whole-tuple accuracy throughout the executed external matrix,
without any claim about the underlying encoder.

\subsection{Why the two metrics disagree}
\label{sec:mechanism}
This guarantee is not cost-free in every scoring view, and
Table~\ref{tab:multilingual} makes the trade-off explicit by counting the
correct field predictions projection creates against those it destroys.
The $N\!/\!A$ net is positive and the
substantive-class net negative in all five fixed arms. If an independently
predicted parent is wrong, projection may replace a correct substantive child
with the structurally required \emph{N/A}; when the parent is correct, the same
operation can repair an unsupported child. Field-wise macro-F1 values these
changes class by class, whereas tuple accuracy credits only four-field
agreement. The field-score effect therefore varies across backbones and
corpora. Post-hoc C-wF1 and hF point in the tuple direction, but were chosen
after the primary analysis and are not needed to establish the pre-specified
secondary result.

Validity alone also does not identify the better structured rule. In Chinese,
M4 scores all 17 legal states jointly yet loses tuple accuracy to greedy M1.
Joint search can improve one field while changing another field that projection
had already made correct. This adverse result cautions against treating a
larger legal-state search as automatically superior; it does not imply that
joint decoding must fail on other models or tasks.

\subsection{Reproducibility}
\label{sec:repro}
Code, pre-registration documents, row-level predictions for all five corpora,
and the pipeline that regenerates every table here are at \url{\repourl}.

\begin{sloppypar}
\paragraph{Checkpoints.} Every backbone is a pinned Hugging Face checkpoint,
its commit revision recorded in the released configuration. Chinese:
\hfid{hfl/chinese-roberta-wwm-ext-large} (anchor),
\hfid{hfl/chinese-roberta-wwm-ext},
\hfid{IDEA-CCNL/Erlangshen-DeBERTa-v2-320M-Chinese},
\hfid{hfl/chinese-electra-180g-large-discriminator}. English:
\hfid{roberta-large}, \hfid{roberta-base},
\hfid{microsoft/deberta-v3-large},
\hfid{google/electra-large-discriminator}. French, Japanese and Korean:
\hfid{FacebookAI/xlm-roberta-large}, \hfid{FacebookAI/xlm-roberta-base},
\hfid{google/rembert}, \hfid{google-bert/bert-base-multilingual-cased}.
\end{sloppypar}

\paragraph{Procedures given only in outline above.} The document-disjoint
split sorts reports by identifier, shuffles them under the seed, re-sorts by
descending row count with a stable sort so the shuffle survives ties, then
assigns each report uniformly at random to one of the two currently lightest
folds, ties broken by fold index. Korean pages come from Poppler's \hfid{pdftotext -f P -l P -layout}, one page per labelled
(URL, page) pair, falling back to \hfid{pdftoppm -r 250} and Tesseract
(\hfid{-l kor+eng}) for image-only pages; the release already carries one
label tuple per unique (URL, page), so no further deduplication applies.
Python 3.10.20, PyTorch 2.2.2+cu121 and Transformers 4.40.2 on Linux; the 32
external arms consumed 25.4 GPU-hours in total on a single RTX 3090.

\subsection{Limitations}
The corpora differ in language, reports, annotation distributions, and primary
backbones, so the replication matrix is not a factorial estimate of a pure
language effect. English and French each have only nine report clusters, and
even Chinese inference resamples just 49. The $32/32$ and $14/32$ summaries
are therefore descriptive directions across executed arms, not a pooled
significance test. AI CUP weights applied to ML-Promise are not an official
ML-Promise metric.

Every contrast here is measured against M0, the unconstrained argmax over our
own probabilities. That is the right comparator for the question we ask, since
it holds rows and probabilities fixed, but the paper therefore reports no
comparison against a published system, an official baseline, or a leaderboard,
and nothing here is a claim about where this system ranks among participants.

Korean has a further input mismatch: its 500 examples are first-384-token PDF
pages rather than released paragraphs, and it has no gold \emph{Misleading}
example. We do not redistribute the extracted Korean text or source PDFs.
The reported Korean predictions are released and can therefore be rescored,
but retraining requires rebuilding the page inputs from their source documents
with the recipe in Section~\ref{sec:repro}.

\begin{sloppypar}
Results for structural loss and other backbones remain exploratory: they show
that reducing the invalid-tuple rate need not yield a conclusive gain on the
official score, but the fixed loss weights and limited checkpoint set cannot
characterize all training objectives or architectures. The rarest Chinese
quality label has two examples and Korean has none, so no class-level claim is
licensed for it. Finally, each seed jointly changes folds and optimization
randomness, so those two sources of variation cannot be separated.
\end{sloppypar}

\section{Conclusion}

Hierarchical projection is not a guaranteed route to higher field-wise
weighted macro-F1. It is instead a retraining-free validity layer that
eliminates invalid tuples by construction and
improves exact whole-tuple accuracy on the Chinese anchor and in all 32
evaluated external arms.
Its field-score effect varies across models and corpora, but its logical guarantee does not.
For audit-oriented tasks with dependent outputs, field-wise performance,
whole-tuple correctness, and structural validity should be reported
together.
